% GENERATED FILE. Edit canonical lessons, then run npm run generate:books.
% canonical-source-hash: fnv1a64:1563123a5db05b89
% canonical-lessons: BN-C465-bibek, BN-C465-bichokkhon, BN-C465-medhabi, BN-C465-kritogyo, BN-C465-bibechona-kora

\chapter{Good Sense}
\label{ch:bn-good-sense}
\hlchaptermodality{465}

\begin{chapteropening}
\textbf{By the end of this chapter:} \emph{I can talk about conscience, good judgement and gratitude.}

Complete the last lesson of chapter 465: I can talk about conscience, good judgement and gratitude.
\end{chapteropening}

\section[\texorpdfstring{\bn{বিবেক} (bibek)}{বিবেক (bibek)}]{\bn{বিবেক} (bibek) — conscience}

\label{lesson:BN-C465-bibek}



\begin{quote}
Before the new one: say the Bengali for indifferent, then the Bengali for to adapt.
\end{quote}

\subsection*{You'll want to know: \bn{বিবেক}}
\textbf{\bn{বিবেক}} — \emph{bibek} — \textquotedblleft{}conscience\textquotedblright{}.

\begin{quote}
\bn{আমার} \bn{বিবেক} \bn{বলছে}, \bn{টাকাটা} \bn{ফিরিয়ে} \bn{দেওয়া} \bn{উচিত।}
\end{quote}

\begin{quote}
\emph{āmār bibek bôlchhe, ṭākāṭā phiriye deoyā uchit.}
\end{quote}

\begin{quote}
My conscience tells me the money ought to be given back.
\end{quote}

\begin{cousinweb}
\textbf{\bn{বিবেক}} (\emph{bibek}) is a tatsama, Sanskrit \emph{viveka}, \textquotedblleft{}discrimination, distinction\textquotedblright{} (Monier-Williams); Bengali uses it for the inner judge that tells right from wrong.
\end{cousinweb}

\begin{grammarlens}[title={What you've built}]
One more word for describing what you lived and felt.
\end{grammarlens}

\subsection*{Your turn}
\begin{itemize}
\raggedright
  \item \emph{Say it:} \emph{bibek}
  \item \emph{Say it:} \emph{bibek}, once more
  \item \emph{Say it:} \emph{āmār bibek bôlchhe, ṭākāṭā phiriye deoyā uchit.} — My conscience tells me the money ought to be given back.
  \item \emph{Recall:} say the Bengali for indifferent, then the Bengali for to adapt, then say \emph{bibek} again
\end{itemize}

\begin{tcolorbox}[breakable,colback=teal!4,colframe=teal!35!black,title={Before you move on}]
What does \bn{বিবেক} mean? (\textquotedblleft{}Conscience\textquotedblright{}.) Say it once more.
\end{tcolorbox}
\section[\texorpdfstring{\bn{বিচক্ষণ} (bichôkkhôn)}{বিচক্ষণ (bichôkkhôn)}]{\bn{বিচক্ষণ} (bichôkkhôn) — shrewd; judicious}

\label{lesson:BN-C465-bichokkhon}



\begin{quote}
Before the new one: say the Bengali for to adapt, then the Bengali for conscience.
\end{quote}

\subsection*{You'll want to know: \bn{বিচক্ষণ}}
\textbf{\bn{বিচক্ষণ}} — \emph{bichôkkhôn} — \textquotedblleft{}shrewd; judicious\textquotedblright{}.

\begin{quote}
\bn{বিচক্ষণ} \bn{ডাক্তার} \bn{রোগটা} \bn{প্রথমেই} \bn{বুঝলেন।}
\end{quote}

\begin{quote}
\emph{bichôkkhôn ḍāktār rogṭā prôthômei bujhlen.}
\end{quote}

\begin{quote}
The shrewd doctor understood the illness straight away.
\end{quote}

\begin{cousinweb}
\textbf{\bn{বিচক্ষণ}} (\emph{bichôkkhôn}) is a tatsama, Sanskrit \emph{vicakṣaṇa}, \textquotedblleft{}conspicuous, clear\textquotedblright{} and \textquotedblleft{}clear-sighted, wise\textquotedblright{} (Monier-Williams).
\end{cousinweb}

\begin{grammarlens}[title={What you've built}]
One more word for describing what you lived and felt.
\end{grammarlens}

\subsection*{Your turn}
\begin{itemize}
\raggedright
  \item \emph{Say it:} \emph{bichôkkhôn}
  \item \emph{Say it:} \emph{bichôkkhôn}, once more
  \item \emph{Say it:} \emph{bichôkkhôn ḍāktār rogṭā prôthômei bujhlen.} — The shrewd doctor understood the illness straight away.
  \item \emph{Recall:} say the Bengali for to adapt, then the Bengali for conscience, then say \emph{bichôkkhôn} again
\end{itemize}

\begin{tcolorbox}[breakable,colback=teal!4,colframe=teal!35!black,title={Before you move on}]
What does \bn{বিচক্ষণ} mean? (\textquotedblleft{}Shrewd; judicious\textquotedblright{}.) Say it once more.
\end{tcolorbox}
\section[\texorpdfstring{\bn{মেধাবী} (medhābī)}{মেধাবী (medhābī)}]{\bn{মেধাবী} (medhābī) — gifted; brilliant (at study)}

\label{lesson:BN-C465-medhabi}



\begin{quote}
Before the new one: say the Bengali for conscience, then the Bengali for judicious.
\end{quote}

\subsection*{You'll want to know: \bn{মেধাবী}}
\textbf{\bn{মেধাবী}} — \emph{medhābī} — \textquotedblleft{}gifted; brilliant (at study)\textquotedblright{}.

\begin{quote}
\bn{মেধাবী} \bn{ছাত্রীটি} \bn{বৃত্তি} \bn{পেল।}
\end{quote}

\begin{quote}
\emph{medhābī chhātrīṭi britti pelô.}
\end{quote}

\begin{quote}
The gifted student won a scholarship.
\end{quote}

\begin{cousinweb}
\textbf{\bn{মেধাবী}} (\emph{medhābī}) is a tatsama, Sanskrit \emph{medhāvin}, \textquotedblleft{}intelligent, wise, a learned man\textquotedblright{} (Monier-Williams).
\end{cousinweb}

\begin{grammarlens}[title={What you've built}]
One more word for describing what you lived and felt.
\end{grammarlens}

\subsection*{Your turn}
\begin{itemize}
\raggedright
  \item \emph{Say it:} \emph{medhābī}
  \item \emph{Say it:} \emph{medhābī}, once more
  \item \emph{Say it:} \emph{medhābī chhātrīṭi britti pelô.} — The gifted student won a scholarship.
  \item \emph{Recall:} say the Bengali for conscience, then the Bengali for judicious, then say \emph{medhābī} again
\end{itemize}

\begin{tcolorbox}[breakable,colback=teal!4,colframe=teal!35!black,title={Before you move on}]
What does \bn{মেধাবী} mean? (\textquotedblleft{}Gifted; brilliant (at study)\textquotedblright{}.) Say it once more.
\end{tcolorbox}
\section[\texorpdfstring{\bn{কৃতজ্ঞ} (kritôgyô)}{কৃতজ্ঞ (kritôgyô)}]{\bn{কৃতজ্ঞ} (kritôgyô) — grateful}

\label{lesson:BN-C465-kritogyo}



\begin{quote}
Before the new one: say the Bengali for judicious, then the Bengali for gifted.
\end{quote}

\subsection*{You'll want to know: \bn{কৃতজ্ঞ}}
\textbf{\bn{কৃতজ্ঞ}} — \emph{kritôgyô} — \textquotedblleft{}grateful\textquotedblright{}.

\begin{quote}
\bn{তোমার} \bn{সাহায্যের} \bn{জন্য} \bn{আমি} \bn{কৃতজ্ঞ।}
\end{quote}

\begin{quote}
\emph{tomār sāhājjer jônyô āmi kritôgyô.}
\end{quote}

\begin{quote}
I am grateful for your help.
\end{quote}

\begin{cousinweb}
\textbf{\bn{কৃতজ্ঞ}} (\emph{kritôgyô}) is a tatsama, Sanskrit \emph{kṛtajña}, \textquotedblleft{}knowing what has been done (for you)\textquotedblright{}, so \textquotedblleft{}grateful\textquotedblright{} (Monier-Williams).
\end{cousinweb}

\begin{grammarlens}[title={What you've built}]
One more word for describing what you lived and felt.
\end{grammarlens}

\subsection*{Your turn}
\begin{itemize}
\raggedright
  \item \emph{Say it:} \emph{kritôgyô}
  \item \emph{Say it:} \emph{kritôgyô}, once more
  \item \emph{Say it:} \emph{tomār sāhājjer jônyô āmi kritôgyô.} — I am grateful for your help.
  \item \emph{Recall:} say the Bengali for judicious, then the Bengali for gifted, then say \emph{kritôgyô} again
\end{itemize}

\begin{tcolorbox}[breakable,colback=teal!4,colframe=teal!35!black,title={Before you move on}]
What does \bn{কৃতজ্ঞ} mean? (\textquotedblleft{}Grateful\textquotedblright{}.) Say it once more.
\end{tcolorbox}
\section[\texorpdfstring{\bn{বিবেচনা} \bn{করা} (bibechônā kôrā)}{বিবেচনা করা (bibechônā kôrā)}]{\bn{বিবেচনা} \bn{করা} (bibechônā kôrā) — to consider; to weigh up}

\label{lesson:BN-C465-bibechona-kora}



\begin{quote}
Before the new one: say the Bengali for gifted, then the Bengali for grateful.
\end{quote}

\subsection*{You'll want to know: \bn{বিবেচনা} \bn{করা}}
\textbf{\bn{বিবেচনা} \bn{করা}} — \emph{bibechônā kôrā} — \textquotedblleft{}to consider; to weigh up\textquotedblright{}.

\begin{quote}
\bn{বাবা} \bn{পরামর্শটা} \bn{বিবেচনা} \bn{করলেন।}
\end{quote}

\begin{quote}
\emph{bābā pôrāmôrshôṭā bibechônā kôrlen.}
\end{quote}

\begin{quote}
Father weighed up the advice.
\end{quote}

Past (he/she …): \bn{তিনি} \bn{সিদ্ধান্তটা} \bn{আবার} \bn{বিবেচনা} \bn{করলেন}, \emph{tini siddhāntôṭā ābār bibechônā kôrlen}.

\begin{cousinweb}
\textbf{\bn{বিবেচনা} \bn{করা}} (\emph{bibechônā kôrā}) is Sanskrit \emph{vivecana}, \textquotedblleft{}discriminating, examining\textquotedblright{} (Monier-Williams), with \emph{kôrā}, \textquotedblleft{}to do\textquotedblright{} (CDIAL 2814).
\end{cousinweb}

\begin{grammarlens}[title={What you've built}]
One more word for describing what you lived and felt.
\end{grammarlens}

\subsection*{Your turn}
\begin{itemize}
\raggedright
  \item \emph{Say it:} \emph{bibechônā kôrā}
  \item \emph{Say it:} \emph{bibechônā kôrā}, once more
  \item \emph{Say it:} \emph{bābā pôrāmôrshôṭā bibechônā kôrlen.} — Father weighed up the advice.
  \item \emph{Recall:} say the Bengali for gifted, then the Bengali for grateful, then say \emph{bibechônā kôrā} again
\end{itemize}

\begin{tcolorbox}[breakable,colback=teal!4,colframe=teal!35!black,title={Before you move on}]
What does \bn{বিবেচনা} \bn{করা} mean? (\textquotedblleft{}To consider; to weigh up\textquotedblright{}.) Say it once more.
\end{tcolorbox}
